\section{Einleitung}

Die geplante Arbeit soll zeigen, wie sich verbale Semantik in Verfassungstexten von 1848 bis 1968 verändert hat. So soll aus rechtslinguistischer Perspektive Evidenz dafür gefunden werden, dass Sprachwandel in Verfassungen mit einem Paradigmenwechsel im Verfassungsrecht korreliert. Die zu beantwortende Forschungsfrage lautet demnach:

\begin{quote}
Ist es möglich, verfassungsrechtliche Paradigmenwechsel durch eine systemisch-funktionale Analyse von Verbphrasen in Verfassungstexten zu identifizieren? Falls ja, wie lassen sich diese grammatisch klassifizieren?
\end{quote}


Somit wird das Thema in das interdisziplinäre Feld der Rechtslinguistik verortet. Dieses ist gerade im Rahmen semantischer Untersuchungen ein wenig erforschter Bereich. Mit dem \textit{konkorp}-Korpus \citep{Thurau2025} liegt nun die Datengrundlage für linguistische Untersuchungen an deutschen Verfassungstexten vor. Das Korpus enthält insgesamt sechs Verfassungstexte aus den Jahren 1848, 1871, 1919, 1949 und 1968. Diese wurden tokenisiert, normalisiert und auf morphosyntaktischer, rechtwissenschaftlich-struktureller und semantischer Ebenen annotiert.

Die Annotationen im \textit{konkorp} ermöglichen es, die zugrundeliegenden Verfassungstexte diachron nach Sprachwandelprozessen zu untersuchen. Besonders von Interesse für die gepalnte Arbeit sind Verbphrasen. Jedoch sollen diese nicht auf formal-syntaktischer Ebene untersucht werden, sondern im Rahmen der \textit{Systemic Functional Linguistics} (SFL, siehe \citealt{Halliday1978,Halliday2004}).

Halliday entwickelte mit der SFL ein Framework, das -- ähnlich wie die Konstruktionsgrammatik (CxG) -- Sprache als ein gebrauchsbasiertes System versteht, welches Bedeutung durch Kontext und Funktionalität darstellt.

SFL basiert im Kontrast zur CxG allerdings weniger auf kognitiven als auf soziolinguistischen Konzepten und stellt Sprache als sozial-funktionales System dar. Gerade für die Untersuchung der verbalen Semantik nutzt die SFL Konzepte zur Satzanalyse, die für die Forschung an Verfassungstexten gut geeignet zu sein scheinen. Dieser Transfer ist Hauptgegenstand der geplanten Arbeit. Das System, dass der Verbphrasenanalyse in der SFL zugrundeliegt, ist das \textit{transitivity system} \citep[170]{Halliday2004}. Anders als in formal-syntaktischen Analysen wird Transitivität hier nicht als eine mögliche Verbvalenz verstanden, sondern als semantisches System. Sätze (in der SFL \textit{clauses}) haben drei funktionale Komponenten; den Vorgang (\textit{process}), die am Vorgang Beteiligten (\textit{participants}) und die Umstände des Vorganges (\textit{circumstances}). Der Vorgang wird durch das Verb im Satz ausgedrückt und kann einer von sechs Kategorien (material, relational, behavioural, metal, verbal oder existential) zugeordnet werden \citep[171]{Halliday2004}.

Das Framework eignet sich - nicht zuletzt aufgrund dieser Kategorisierung -, um es auf schriftsprachliche Daten anzuwenden. Gleichzeitig ist die interdisziplinäre Forschung von Verfassungstexten, gerade im Bezug auf Verbsemantik, noch nicht hinreichend erforscht. \cite{Busse2015,BusseEtAl2018} führten Frame-semantische Studien an Gesetzestexten durch und untersuchten dabei Nomina aus dem Zivil- und Strafrecht, \cite{Cirko2022} forschte an Nomina in Grundrechten des deutschen Grundgesetzes und der polnischen Verfassung, und \cite{Sayatz1991} schrieb einen Artikel zur illokutiven Interpretation von Deklarativsätzen mit Modalverben - jedoch findet sich keine korpusbasierte Studie zur diachroner Rechtslinguistik in Bezug auf Verben in Verfassungstexten. Dabei ist die Relevanz grammatischer Analysen für die Rechtswissenschaft nicht abzustreiten. Seit den 1980ern untersucht die sogennante Heidelberger Gruppe rechtslinguistische Themen, und auch international ist die Disziplin der \textit{Legal linguistics} in der Sprachwissenschaft gut verankert. Demnach soll die geplante Arbeit auch dazu dienen, die oben beschriebene Forschungslücke zu füllen. Durch einen systemisch-funktionalen Ansatz soll ein Framework genutzt werden, das bis dato für die Rechtslinguistik noch nicht verwendet worden ist. Sollte es gelingen, SFL auf die Rechtslinguistik zu übertragen, kann dies ein Schritt Richtung methodologischer Vielfalt zur Erschließung der Semantik von Gesetzestexten sein.

%Dabei könnte gerade eine funktionale semantische Analyse Einsicht darin bieten, wie sich die Verhältnisse von Staat und Bürger*innen diachron verändert haben.


%welche Relationen zwischen Staat und Bürger*innen vorlagen, wer auf wen in welcher Weise normativ einwirkte und zu wessen Gunsten welche Normen erlassen worden sind.


%Diese und weitere Aspekte sollen in der geplanten Arbeit untersucht werden.

\section{Theoretischer Hintergrund}\label{sec:theory}

Für die theoretischen Grundlagen der Arbeit sollen die notwendigen theoretischen Konzepte eingeführt und ihr Nutzen für die Arbeit begründet werden. Diese sind:

\begin{enumerate}
\item \textbf{Systemic Functional Linguistics (SFL).}\\ Zunächst soll das Framework der SFL nach Michael Halliday und die Relevanz der vier Ebenen, auf denen sprachliche Äußerungen (auch in Textform) analysiert werden können, vorgestellt werden. Diese Ebenen sind: \textit{Phonology}, \textit{Context}, \textit{Lexico-grammar} und \textit{Semantics}. Bis auf die Phonologie sollen in die geplante Analyse alle Ebenen mit einbezogen werden. Ebenso soll kurz erläutert werden, warum die Datengrundlage Anreize zum Arbeiten in diesem Framework schafft.

Besonderer Fokus wird auf das Konzept der Transitivität gelegt. Anders als in der formalen Syntak wird Transitivität in der SFL als eine funktionale Kategorie zur Erschließung von Satzbedeutungen betrachtet. Diese setzt sich zusammen aus a) den Vorgängen, die beschrieben werden, b) den an den Vorgängen Beteiligten und c) den Umständen der Vorgänge. Später wird in der Methodik beschrieben, wie dieses Konzept zur qualitativen Analyse von Verbclustern in der vorliegenden Datengrundlage genutzt wird.
%\item \textbf{Register.} \\
%Da die Kontext-Ebene in der SFL eng verwoben mit dem Registerbegriff ist, wird dieser als nächstes eingeführt. An dieser Stelle soll kurz darauf verwiesen werden, dass der Registerbegriff für Schriftsprache ggf. anders konnotiert werden muss als für gesprochene Sprache. Sollte sich im Verlauf der Arbeit herausstellen. dass der Transfer Herausforderungen mit sich bringt, wird auf eine ausführlichere Diskussion im Analyseabschnitt verwiesen.
\item \textbf{Rechtslinguistik.}\\ In diesem Unterabschnitt wird ein Überblick über die Forschungsschwerpunkte dieser linguistischen Disziplin gegeben. Anders als in \citealt{Thurau2025} ist geplant, Rechtslinguistik auch kurz von der forensischen Linguistik abzugrenzen. Jedoch möchte ich, ähnlich wie \citeauthor{Thurau2025}, durch die Arbeit betonen, dass Rechtssprache ein gänzlich eigenes Register skizziert.
\item \textbf{Rechtsgeschichte.}\\ Für den nötigen interdisziplinären Kontext folgt eine kurze Darstellung der Disziplin er Rechtsgeschichte \citep{Hähnchen2016}. Um gerade das rechtshistorische Fundament für die Datengrundlage der Studie zu schaffen, wird ein Fokus auf das deutsche Verfassungsrecht gelegt \citep{WeberFas1983,IpsenEtAl2022}.
\item \textbf{Korpuslinguistik und das \textit{konkorp}.}\\ Anschließend wird die Disziplin der Korpuslinguistik anhand der Datengrundlage für diese Arbeit vorgestellt und die Vorteile von multi-layered Korpora erläutert \citep{OdebrechtEtAl2016}.  Das zu untersuchende Korpus ist hierbei das \textit{konkorp}. Speziell werden Studien aus \citealt{VogelEtAl2022} genutzt, um den Nutzen interdisziplinärer Forschung zwischen der Jurisprudenz und Korpuslinguistik darzustellen.
\end{enumerate}

\noindent Ziel ist es, die genannten Konzepte in der qualitativen Analyse so miteinander zu verknüpfen, dass die oben genannte Forschungsfrage beantwortet werden kann. Im Fokus wird dabei der methodologische Versuch des Transfers der SFL auf die Korpuslinguistik am Beispiel des \textit{konkorp} sein.

\section{Methodik}

In diesem Abschnitt soll das konkrete Vorgehen zur Beantwortung der Forschungsfrage dargestellt werden. Er wird in einen quantitativen und einen qualitativen Teil gegliedert.

Zunächst wird für ersteren das \textit{konkorp} in ANNIS eingelesen und mittels regulärer Ausdrücke nach (nach syntaktischer Definition) Verben durchsucht. Es wird eine Frequenzanalyse von Verbphrasen durchgeführt. Dies geschieht für jeden der sechs Datensätze separat, um die Ergebnisse auch diachron auswerten zu können.

Es werden schließlich qualitative Analysen einzelner Verbokkurenzen vorgenommen. Hierzu wird auf die Dreiteilung des Transitivitätskonzeptes (processes, participants, circumstances) nach Halliday zurückgegriffen.
Halliday schreibt Sätzen (\textit{clauses}) drei Metaebenen von Bedeutung zu; Austausch, Nachricht, Repräsentation. Die Analyse wird sich vorranging auf die letzte Ebene beziehen, denn auf dieser findet die Konzeption von Transitivität in Hallidays Framework statt.

Auf strukturalistischer Ebene korrelieren Komponenten von Transititivät mit syntaktischen Einheiten (process -- Verbphrase, participants -- Nominalphrase, circumstances -- Präpositionalphrase und Adverbien). Da die Annotation von Konstituenten im \textit{konkorp} bereits erfolgt ist, muss hier keine Analyse mehr erfolgen.
Vielmehr wird das Korpus um zwei Annotationsebenen erweitert; Verbkomplexität und \textit{process type}. In ersterer wird annotiert, ob es sich um ein Verbum simplex, ein Präfix- oder ein Partikelverb handelt. In zweiterer werden die Verben einer der sechs Kategorien nach \citet[171]{Halliday2004} zugeordnet und jeweils die Umstände und Beteiligten des durch das Verb beschriebenen Vorganges identifiziert. Dies soll für alle sechs zugrundeliegenden Verfassungen geschehen. So sollen die einzelnen Verfassungstexte anhand der Nutzung der jeweiligen Verben charakterisiert werden. Ziel ist es, die Hypothese zu überprüfen, ob die Verfassungen mit zunehmender Aktualität einen Wandel in der Dynamik zwischen Staat und Bürger*innen verzeichnen lassen. Hierzu wird in der Analyse und Diskussion betrachtet, welche \textit{process types} in welchem Datensatz wie häufig vorhanden sind und ob sich durch die jeweils in der \textit{clause} enthaltenen \textit{participants} und \textit{circumstances} Rückschlüsse auf verfassungsrechtliche Konzepte ziehen lassen.

So könnte beispielsweise untersucht werden, welche Rolle den Staatsbürger*innen und welche dem Staat zugeordnet werden können. Welches Organ hat in welcher Epoche welche Rechte und Pflichten, wer handelt, und wer wird behandelt?

Im Idealfall kann linguistische Evidenz dafür gefunden werden, dass aktuellere Verfassungstexte dem Schutz und dem Einräumen von Rechten gegenüber den Bürger*innen dienen, während ältere Verfassungen das Ziel hatten, die Staatsgewalt und ihre Macht zu sichern.


%\subsection{Exemplarische Analyse eines transitiven Verbs}

%clause as exchange, as messange, as representation

%auf ebene word bezogen
%Paulskirchverfassung 1849
%Verfassuung des deutschen Reiches 1871
% https://www.verfassungen.de/de67-18/verfassung71-i.htm
%die Verfassung des deutschen Reiches 1919
%Grundgesetz für die Bundesrepublik Deutschland 1949
%Die Verfassung der Deutschen Demokratischen Republik 1949
%Verfassung der Deutschen Demokratischen Republik 1968


% pos=/V.*/ -> 6057
